% ================================================================
% PAGE 4 — PROBLEM, RESEARCH QUESTIONS, OBJECTIVES, SCOPE
% ================================================================

\section{Problem}

\begin{itemize}
  \item Thresholds are computed after federated training.
  \item Thresholds define the final anomaly decision boundary.
  \item Benign calibration scores determine~$\tau$.
  \item Under non-IID device heterogeneity, one shared~$\tau$ may not fit all clients.
  \item Per-client FPR disparity is operationally important because false alarms are experienced locally.
\end{itemize}

\subsection{Central Research Question}

\begin{tcolorbox}[callout]
\itshape
How much does threshold calibration scope reduce per-client false-positive-rate disparity
under a fixed FedAvg autoencoder in non-IID federated IoT anomaly detection?
\end{tcolorbox}

\subsection{Research Questions}

\begin{itemize}
  \item \textbf{RQ1:} How large is per-client FPR disparity under the client-averaged shared threshold?
  \item \textbf{RQ2:} Does per-client threshold calibration reduce CV(FPR) relative to the
    client-averaged shared threshold under Regime A?
  \item \textbf{RQ3:} Can cluster-mean calibration approximate per-client calibration
    without a device taxonomy?
\end{itemize}

\section{Objectives}

\begin{minipage}[t]{0.48\linewidth}
\subsection*{Scientific objectives}
\begin{itemize}
  \item Isolate threshold calibration scope as the experimental variable.
  \item Quantify per-client FPR dispersion under B1.
  \item Measure B1-to-B2 CV(FPR) reduction.
  \item Evaluate B4 as a taxonomy-free approximation.
  \item Characterize detection-quality tradeoffs.
\end{itemize}
\end{minipage}%
\hfill
\begin{minipage}[t]{0.48\linewidth}
\subsection*{Technical objectives}
\begin{itemize}
  \item Preserve fixed FedAvg autoencoder training.
  \item Reuse the same score artifacts across threshold policies.
  \item Apply eligibility rules consistently.
  \item Report CV(FPR), CV(TPR), worst-client BA, Macro-F1, and P10 Macro-F1.
  \item Preserve reproducibility of tables, figures, and seeds.
\end{itemize}
\end{minipage}

\vspace{0.5em}

\section{Scope}

\vspace{0.2em}
\begin{tabularx}{\linewidth}{@{} >{\small}X >{\small}X @{}}
\toprule
\textbf{Included} & \textbf{Excluded}\\
\midrule
N-BaIoT Regime A as primary study        & Training-data poisoning\\[2pt]
Physical device clients                  & Model poisoning\\[2pt]
Fixed FedAvg autoencoder                 & Aggregation poisoning\\[2pt]
B1 client-averaged shared threshold      & Backdoors\\[2pt]
B2 per-client threshold                  & Evasion\\[2pt]
B3 family-mean threshold, Regime A only  & Privacy mechanisms\\[2pt]
B4 cluster-mean threshold                & Deployment benchmarking\\[2pt]
CICIoT2023 file-level applicability boundary & Hardware validation\\[2pt]
Regime C Dirichlet severity sweep        & Concept drift\\[2pt]
CV(FPR) over eligible clients            & Model-personalization comparators\\[2pt]
B0 centralized reference comparator      & Universal IoT malware detection claims\\[2pt]
                                          & Formal privacy guarantees\\
\bottomrule
\end{tabularx}

\clearpage
